\documentclass[11pt]{article}
\usepackage[a4paper,margin=18mm]{geometry}
\usepackage[no-math]{fontspec}
\setmainfont{Latin Modern Roman}
\usepackage{amsmath,amssymb,amsthm,xfrac,xspace,hyperref,graphicx,comment}
\excludecomment{DefTralics}
\providecommand{\C}{}
\providecommand{\bibliofont}{\small}
\newcommand{\ie}{i.e.,\xspace}
\newcommand{\eg}{e.g.,\xspace}
\newcommand{\etc}{etc.\xspace}
\newcommand{\cf}{cf.\xspace}
\newcommand{\resp}{resp.\xspace}
\newcommand{\smaller}{\small}
\newcommand{\Subsection}{\subsection}
\newcommand{\Subsubsection}{\subsubsection}
\newcommand{\skpt}{\smallskip}
\emergencystretch=3em
\newcommand{\OneNumberAllTheorems}{}
\theoremstyle{plain}
\newtheorem{theorem}{Theorem}[section]
\newtheorem{proposition}[theorem]{Proposition}
\newtheorem{lemma}[theorem]{Lemma}
\newtheorem{corollary}[theorem]{Corollary}
\theoremstyle{definition}
\newtheorem{definition}[theorem]{Definition}
\newtheorem{remark}[theorem]{Remark}
\newtheorem{example}[theorem]{Example}
\newtheorem*{namedassumption}{\jepassumptionname}
\newenvironment{enonce*}[1]{\def\jepassumptionname{#1}\begin{namedassumption}}{\end{namedassumption}}
\input{author-preamble.tex}
\begin{document}
\begin{center}{\Large Stable item 1934}\end{center}
\paragraph{Source and attribution.}
Sergei Kuksin, Vahagn Nersesyan, Armen Shirikyan, \emph{Mixing via controllability for randomly forced nonlinear dissipative PDEs}, Journal de l’École polytechnique — Mathématiques 7 (2020), publisher version, DOI 10.5802/jep.130. \href{https://jep.centre-mersenne.org/articles/10.5802/jep.130/}{Publisher article}. Authors' text: \href{https://creativecommons.org/licenses/by/4.0/}{CC BY 4.0}.
\paragraph{Editorial problem boundary.}
Prove only Theorem 3.1 below: under the original stabilisation condition (A) and measure-transformation condition (B), the displayed transition probabilities are uniformly continuous in the dual-Lipschitz distance over all iterates. The stopped triangular path transformation, full total-variation estimate, and complete supporting Lemma 3.3 proof are retained. The selected statement assumes (A) and (B); stationary-measure, mixing and PDE applications are outside this item.
\paragraph{Difficulty (editorial): H2.}
Stopping the triangular noise transformation at a distance barrier yields a summable total-variation defect on the infinite product noise law. The transformed trajectories agree on a high-probability contracting event, converting approximate rather than exact coupling into equicontinuity uniformly over all Markov iterates.
\paragraph{Editorial transcription note.}
Original author language, mathematics and ordering are preserved. Publisher front matter is replaced by this independent article wrapper. Bibliography is recovered from matching cached publisher HTML, including its TeX formulas. No new proof or mathematical repair is supplied. Surrounding original results are retained for conservative context and are not extra selected corpus claims.
\clearpage
\input{author-body.tex}
\begin{thebibliography}{99}
\bibitem{AS-2005} A. A. Agrachev \& A. V. Sarychev - “Navier–Stokes equations: controllability by means of low modes forcing” , J. Math. Fluid Mech. 7 (2005) no. 1, p. 108-152
\bibitem{AS-2006} A. A. Agrachev \& A. V. Sarychev - “Controllability of 2D Euler and Navier–Stokes equations by degenerate forcing” , Comm. Math. Phys. 265 (2006) no. 3, p. 673-697
\bibitem{BKL-2002} J. Bricmont, A. Kupiainen \& R. Lefevere - “Exponential mixing of the 2D stochastic Navier–Stokes dynamics” , Comm. Math. Phys. 230 (2002) no. 1, p. 87-132
\bibitem{bricmont-2002} J. Bricmont - “Ergodicity and mixing for stochastic partial differential equations” , in Proceedings of the International Congress of Mathematicians, Vol. I (Beijing, 2002) , Higher Ed. Press, Beijing, 2002, p. 567-585
\bibitem{BV1992} A. V. Babin \& M. I. Vishik - Attractors of evolution equations , North-Holland Publishing, Amsterdam, 1992
\bibitem{CI-1974} N. Chafee \& E. F. Infante - “A bifurcation problem for a nonlinear partial differential equation of parabolic type” , Applicable Anal. 4 (1974), p. 17-37
\bibitem{debussche-2013} A. Debussche - “Ergodicity results for the stochastic Navier–Stokes equations: an introduction” , in Topics in mathematical fluid mechanics , Lect. Notes in Math. , vol. 2073 , Springer, Heidelberg, 2013, p. 23-108
\bibitem{EMS-2001} W. E, J. C. Mattingly \& Y. Sinai - “Gibbsian dynamics and ergodicity for the stochastically forced Navier–Stokes equation” , Comm. Math. Phys. 224 (2001) no. 1, p. 83-106
\bibitem{ES-2000} W. E \& Y. G. Sinaĭ - “New results in mathematical and statistical hydrodynamics” , Russian Math. Surveys 55 (2000) no. 4, p. 635-666
\bibitem{feller1968} W. Feller - An introduction to probability theory and its applications. Vol. I , John Wiley \& Sons, New York, 1968
\bibitem{FGRT-2015} J. Földes, N. Glatt-Holtz, G. Richards \& E. Thomann - “Ergodic and mixing properties of the Boussinesq equations with a degenerate random forcing” , J. Funct. Anal. 269 (2015) no. 8, p. 2427-2504
\bibitem{FM-1995} F. Flandoli \& B. Maslowski - “Ergodicity of the 2D Navier–Stokes equation under random perturbations” , Comm. Math. Phys. 172 (1995) no. 1, p. 119-141
\bibitem{FW1984} M. I. Freidlin \& A. D. Wentzell - Random perturbations of dynamical systems , Springer, New York–Berlin, 1984
\bibitem{HM-2006} M. Hairer \& J. C. Mattingly - “Ergodicity of the 2D Navier–Stokes equations with degenerate stochastic forcing” , Ann. of Math. (2) 164 (2006) no. 3, p. 993-1032
\bibitem{HM-2011} M. Hairer \& J. C. Mattingly - “A theory of hypoellipticity and unique ergodicity for semilinear stochastic PDEs” , Electron. J. Probab. 16 (2011) no. 23, p. 658-738
\bibitem{JNPS-cmp2014} V. Jakšić, V. Nersesyan, C.-A. Pillet \& A. Shirikyan - “Large deviations and Gallavotti–Cohen principle for dissipative PDE’s with rough noise” , Comm. Math. Phys. 336 (2015) no. 1, p. 131-170
\bibitem{KNS-2020} S. Kuksin, V. Nersesyan \& A. Shirikyan - “Exponential mixing for a class of dissipative PDEs with bounded degenerate noise” , Geom. Funct. Anal. 30 (2020) no. 1, p. 126-187
\bibitem{KS-cmp2000} S. Kuksin \& A. Shirikyan - “Stochastic dissipative PDEs and Gibbs measures” , Comm. Math. Phys. 213 (2000) no. 2, p. 291-330
\bibitem{KS-book} S. Kuksin \& A. Shirikyan - Mathematics of two-dimensional turbulence , Cambridge Tracts in Math. , vol. 194 , Cambridge University Press, Cambridge, 2012
\bibitem{kuksin-1982} S. Kuksin - “Diffeomorphisms of function spaces that correspond to quasilinear parabolic equations” , Mat. Sb. (N.S.) 117(159) (1982) no. 3, p. 359-378, 431
\bibitem{lamperti1996} J. W. Lamperti - Probability , John Wiley \& Sons, New York, 1996
\bibitem{martirosyan-2017} D. Martirosyan - “Large deviations for stationary measures of stochastic nonlinear wave equations with smooth white noise” , Comm. Pure Appl. Math. 70 (2017) no. 9, p. 1754-1797
\bibitem{nersisyan-2010} H. Nersisyan - “Controllability of 3D incompressible Euler equations by a finite-dimensional external force” , ESAIM Control Optim. Calc. Var. 16 (2010) no. 3, p. 677-694
\bibitem{nersesyan-2020} V. Nersesyan - “Approximate controllability of nonlinear parabolic PDEs in arbitrary space dimension” , Math. Control Relat. Fields (2020), under revision
\bibitem{shirikyan-cmp2006} A. Shirikyan - “Approximate controllability of three-dimensional Navier–Stokes equations” , Comm. Math. Phys. 266 (2006) no. 1, p. 123-151
\bibitem{shirikyan-asens2015} A. Shirikyan - “Control and mixing for 2D Navier–Stokes equations with space-time localised noise” , Ann. Sci. École Norm. Sup. (4) 48 (2015) no. 2, p. 253-280
\bibitem{shirikyan-2020} A. Shirikyan - “Controllability implies mixing II. Convergence in the dual-Lipschitz metric” , J. Eur. Math. Soc. (JEMS) (2020), accepted for publication
\bibitem{smale-1965} S. Smale - “An infinite dimensional version of Sard’s theorem” , Amer. J. Math. 87 (1965), p. 861-866
\bibitem{zabczyk2008} J. Zabczyk - Mathematical control theory , Modern Birkhäuser Classics , Birkhäuser, Boston, MA, 2008
\end{thebibliography}
\end{document}
